% 자동 생성 --- tools/make_tables_15p.py (지표 9개, 돌파 수 제외, 거리는 ell 단위). 손으로 고치지 말 것.
\begin{table*}[t]
\centering
\caption{2$\times$2 상호작용 분해 (Gemini 지휘관). 9개 지표 중 \textbf{8개}에서 상호작용이 유의하지 않으므로 두 계층의 기여는 \textbf{가산적}이다.}
\label{tab:interaction}
\footnotesize
\resizebox{\ifdim\width>\linewidth\linewidth\else\width\fi}{!}{%
\begin{tabular}{lcc}
\toprule
지표 & 총효과 & 상호작용 [95\,\% CI] \\
\midrule
포획률 & +0.062 & -0.039 [-0.083, +0.012] \\
충돌률(척당) & -0.052 & -0.015 [-0.070, +0.022] \\
그물 걸림 & -0.60 & +0.17 [-0.08, +0.41] \\
포획당 그물 & -0.328 & +0.004 [-0.045, +0.059] \\
포획당 이동거리 ($\ell$) & -0.72 & -0.22 [-0.50, +0.12] \\
총 선회량 (rad) & -757 & -56 [-165, +59] \\
포획 이격거리 ($\ell$) & +0.88 & +0.15 [-0.01, +0.32] \\
평균 포획 시각 (step) & -22 & -12 [-20, -6]$^{*}$ \\
전개 시 적 거리 ($\ell$) & -1.47 & +0.10 [-0.17, +0.43] \\
\bottomrule
\end{tabular}
}
\\[3pt]
\begin{minipage}{\linewidth}\footnotesize
상호작용 $=$ (LLM+U-Net $-$ LLM+휴리스틱) $-$ (휴리스틱+U-Net $-$ baseline). 표식 $^{*}$ 은 유의한 상호작용이며, 나머지는 두 계층이 서로를 돕지도 잡아먹지도 않는다는 뜻이다.
\end{minipage}
\end{table*}
